\setlength{\baselineskip}{20pt}
\begin{englishabstract}

Diffusion models, as a class of powerful generative frameworks, have demonstrated outstanding performance in fields such as image synthesis and audio generation by progressively adding Gaussian noise and training denoising networks for reverse prediction. However, the potential of diffusion models in discriminative tasks remains largely underexplored. Representation learning, as the core component of discriminative visual tasks, aims to learn highly discriminative feature representations for downstream applications including image classification, object detection, and image segmentation. Currently, a few studies have attempted to introduce diffusion mechanisms into specific discriminative scenarios, but these methods are typically tailored to individual tasks and lack generality. Moreover, incorporating recursive denoising operations during inference introduces significant additional computational overhead. Therefore, how to integrate the denoising capability of diffusion models into discriminative representation learning in a general and efficient manner, enhancing feature discriminability without increasing inference cost, remains a critical open problem.

This thesis conducts a systematic study on feature enhancement algorithms for image discriminative tasks based on diffusion models. The main contributions are as follows:

(1) Research on computation-free feature enhancement algorithms. This work innovatively reinterprets each embedding layer in the backbone network as a denoising layer, drawing an analogy between the $N$ cascaded embedding layers and the $T$-step recursive denoising process in diffusion models, achieving unified modeling of feature extraction and feature denoising within a single framework. Building upon this, a parameter fusion strategy is derived through rigorous mathematical proof, demonstrating that the denoising layer parameters can be equivalently merged into the original embedding layers, enabling zero additional computational overhead for feature denoising during inference. This method operates in a label-free manner to improve feature quality, while also serving as a complementary mechanism to further enhance performance when labels are available. Experimental results on multiple visual tasks, including person re-identification, image classification, object detection, and image segmentation, demonstrate that the proposed method can stably and significantly improve the performance of various baseline models.

(2) Research on distillation-based feature enhancement algorithms. To address the limited expressiveness of lightweight denoising modules and the insufficient denoising effectiveness under constrained depth, two synergistic enhancement strategies are further proposed. On one hand, a teacher--student distillation framework is introduced, leveraging a high-capacity teacher network to provide more accurate noise distribution estimation guidance for the lightweight student denoising module, effectively calibrating the denoising direction and reducing prediction bias, with the additional overhead incurred only during training. On the other hand, a multi-step denoising fusion strategy is proposed, comprising multi-round single-layer denoising and progressive distillation-based multi-step fusion as two complementary mechanisms, which mathematically collapse multi-step iterative processes into equivalent single-step linear operations, significantly enhancing denoising capability under limited layers. The two strategies complement each other from the perspectives of noise prediction accuracy and effective denoising amplitude, jointly achieving superior feature discriminability. Extensive experiments validate the effectiveness and generalization ability of the proposed methods across multiple backbone architectures, while consistently maintaining the core advantage of zero additional computational overhead during inference.

\par\noindent\englishkeywords{Diffusion Model; Representation Learning; Feature Denoising; Parameter Fusion; Knowledge Distillation}
\end{englishabstract}